\documentclass[12pt]{article}
\usepackage{amsmath,amssymb,amsthm}

\begin{document}

\begin{proof}
Proof name: Cea's lemma (Galerkin error estimate).

Let $y \in H^1_0(\Omega)$ be the solution of the continuous problem
\[
a(y,v)=F(v) \quad \text{for all } v\in H^1_0(\Omega),
\]
and let $y_h \in V^0_h$ be the solution of the discrete problem
\[
a(y_h,v_h)=F(v_h) \quad \text{for all } v_h\in V^0_h,
\]
where $a(\cdot,\cdot)$ is bounded and coercive on $H^1_0(\Omega)$ with constants $C$ and $c$, respectively:
\[
|a(u,v)| \le C \|u\|_{1,\Omega} \|v\|_{1,\Omega}
\]
and
\[
a(v,v) \ge c \|v\|_{1,\Omega}^2
\]
for all $v\in H^1_0(\Omega)$. By Lax--Milgram, such $y$ and $y_h$ exist uniquely.

Define the error $e := y - y_h$. Then Galerkin orthogonality yields
\[
a(e, v_h)=0 \quad \text{for all } v_h\in V^0_h.
\]
For any $w_h\in V^0_h$, decompose $e$ as
\[
e = (y - w_h) + (w_h - y_h).
\]
Since $a(e, w_h - y_h)=0$ by orthogonality, we obtain
\[
a(e,e)=a(e, y - w_h).
\]

Using the boundedness of $a$, we have
\[
|a(e,e)| \le C \|e\|_{1,\Omega} \|y - w_h\|_{1,\Omega}.
\]
Using coercivity,
\[
a(e,e) \ge c \|e\|_{1,\Omega}^2.
\]
Therefore
\[
c \|e\|_{1,\Omega}^2 \le C \|e\|_{1,\Omega} \|y - w_h\|_{1,\Omega}.
\]
If $e \equiv 0$ the desired inequality holds, otherwise we may divide by $\|e\|_{1,\Omega}$ to obtain
\[
\|e\|_{1,\Omega} \le (C/c) \|y - w_h\|_{1,\Omega}.
\]
This holds for every $w_h \in V^0_h$, hence taking the infimum over $w_h \in V^0_h$ gives
\[
\|y - y_h\|_{1,\Omega} = \|e\|_{1,\Omega} \le (C/c) \inf_{w_h\in V^0_h} \|y - w_h\|_{1,\Omega}.
\]
This is precisely the Cea inequality: the Galerkin error is controlled by the best approximation error in the energy norm, up to the constant $C/c$ depending only on the boundedness and coercivity constants of $a$.
\end{proof}

\end{document}